%! Elimination Matrices and Inverse Matrices — Unit 2, Lesson 2.2 — Warm-Up (Answer Key)
\documentclass[10pt]{article}
\usepackage{linalg-article}
\usepackage{linalg-key}   % pulls in -boxes; do NOT also load -boxes
\newcommand{\TallMath}[1]{$\displaystyle #1\rule[-1.4em]{0pt}{3.2em}$}
\newcommand{\xx}{\mathbf{x}}
\newcommand{\bb}{\mathbf{b}}

\begin{document}

\pageheader{Unit 2, Lesson 2.2}{Warm-Up --- Answer Key}
\namedateperiod

\vspace{0.12in}

\begin{notesbox}{Getting Ready: Skills We'll Use Today}
\small Today every elimination step becomes a \emph{matrix}. These three moves are the pieces.
Answer quickly.

\vspace{4pt}
\textbf{1. Matrix times a vector} (Lesson 1.3). Multiply row by row:
\[
  \begin{bmatrix} 1 & 0 \\ -2 & 1 \end{bmatrix}
  \begin{bmatrix} 5 \\ 12 \end{bmatrix}
  = \begin{bmatrix}{\color{keyred}\mathbf{5}}\\[2pt]{\color{keyred}\mathbf{2}}\end{bmatrix}.
\]
\hspace{2em}{\footnotesize(top: $1\cdot5+0\cdot12=5$; bottom: $-2\cdot5+1\cdot12=2$)}

\vspace{4pt}
\textbf{2. Matrix times a matrix} (Lesson 1.4). Take the product one column at a time:
\[
  \begin{bmatrix} 1 & 0 \\ -3 & 1 \end{bmatrix}
  \begin{bmatrix} 2 & 1 \\ 6 & 5 \end{bmatrix}
  = \begin{bmatrix}{\color{keyred}\mathbf{2}} & {\color{keyred}\mathbf{1}}\\[2pt]{\color{keyred}\mathbf{0}} & {\color{keyred}\mathbf{2}}\end{bmatrix}.
\]
Notice the top row is unchanged, and the bottom row became (row~2)$\,-3\,$(row~1).
\hspace{1em}{\footnotesize($-3(2,1)+(6,5)=(0,2)$)}

\vspace{4pt}
\textbf{3. One multiplier} (Lesson 2.1). To clear the $x$ below the pivot in
$\ x + 3y = 7,\ \ 2x + 7y = 15$, the pivot is $1$ and the multiplier is
\[
  \ell = \frac{\text{entry to eliminate}}{\text{pivot}} = \frac{2}{1} = {\color{keyred}\mathbf{2}}.
\]
\end{notesbox}

\begin{teachernote}
Items~1--2 are the exact shapes students meet in the notes: $E\bb$ and $EA$. Item~2 previews the
punch line --- multiplying by $\begin{bsmallmatrix} 1&0\\-\ell&1 \end{bsmallmatrix}$ leaves row~1
alone and performs one elimination step on row~2. Item~3's system and multiplier ($A=\begin{bsmallmatrix} 1&3\\2&7 \end{bsmallmatrix}$, $\ell=2$)
are the running example carried through the notes.
\end{teachernote}

\end{document}
